\documentclass[11pt]{article}
\usepackage[margin=1in]{geometry}
\usepackage{amsmath,amssymb}
\usepackage{xcolor}
\usepackage{booktabs}
\usepackage{enumitem}
\usepackage{fancyhdr}
\usepackage{tcolorbox}
\tcbuselibrary{skins,breakable}
\usepackage{hyperref}
\hypersetup{colorlinks=true, linkcolor=blue, urlcolor=blue}

\definecolor{primary}{RGB}{30,64,140}
\definecolor{tipbg}{RGB}{235,244,255}
\definecolor{examplebg}{RGB}{240,248,235}
\definecolor{formulabg}{RGB}{255,247,224}
\definecolor{warnbg}{RGB}{255,236,236}

\pagestyle{fancy}
\fancyhf{}
\renewcommand{\headrulewidth}{0.4pt}
\cfoot{\thepage}

\newtcolorbox{tipbox}[1][Tip]{
  colback=tipbg, colframe=primary, fonttitle=\bfseries,
  title=#1, breakable, rounded corners
}
\newtcolorbox{examplebox}[1][Example]{
  colback=examplebg, colframe=primary!70!black, fonttitle=\bfseries,
  title=#1, breakable, rounded corners
}
\newtcolorbox{formulabox}[1][Formula]{
  colback=formulabg, colframe=primary!70!black, fonttitle=\bfseries,
  title=#1, breakable, rounded corners
}
\newtcolorbox{warnbox}[1][Warning]{
  colback=warnbg, colframe=red!60!black, fonttitle=\bfseries,
  title=#1, breakable, rounded corners
}


\title{\color{primary}\textbf{Time \& Work --- Remaining Variants}\\[4pt]\large The Last 2--5\%: Rarer TCS Digital Phrasings + Extra Trigger-Word Tricks}
\author{Aptitude Grind}
\date{}

\lhead{\small Time \& Work --- Remaining Variants}
\rhead{\small Aptitude Grind}

\begin{document}
\maketitle
\thispagestyle{fancy}

\noindent Your first three booklets (\texttt{01-fundamentals}, \texttt{02-fast-tricks-and-shortcuts},
\texttt{03-question-pattern-catalog}) cover essentially all high-frequency TCS Ninja + Digital Time \&
Work questions. This file adds the small remainder: ten less-common variants and seven trigger-word
tricks that occasionally show up, mostly on the Digital-level paper. Treat this as a short top-up,
not a new foundation --- every method below still reduces to the core ideas in
\texttt{01-fundamentals.tex}.

\section*{Category K --- Rare But Seen Variants}

\subsection*{K1. Working for x Days, Then Leaving}
\begin{tipbox}[Phrasing]
``A can finish a job in 20 days. A works for 6 days and then leaves. B finishes the remaining work
alone in 21 days. In how many days can B alone finish the whole job?''
\end{tipbox}
\textbf{Method:} Completed work $\to$ remaining work $\to$ solve. A's 6-day share $=6/20=3/10$;
remaining $=7/10$. If B does that $7/10$ in 21 days, B's full-job time $=21/(7/10)=30$ \textbf{days}.

\subsection*{K2. Working Together, Then One Leaves}
\begin{tipbox}[Phrasing]
``A and B together can finish a job in 8 days. They work together for 5 days, then A leaves. B
finishes the rest alone in 6 more days. In how many days can B alone do the whole job?''
\end{tipbox}
\textbf{Method:} Joint work $\to$ remaining work $\to$ single worker. Combined rate $=1/8$; 5-day
joint share $=5/8$; remaining $=3/8$. B does $3/8$ in 6 days $\Rightarrow$ B's full-job time
$=6/(3/8)=16$ \textbf{days}.

\subsection*{K3. Fraction of Work Completed by Each Person}
\begin{tipbox}[Phrasing]
``A completed 3/5 of a work in 12 days, then left. B completed the remaining work in 6 days. Find the
time each would take to complete the full work alone, and the time taken if they had worked together
from the start.''
\end{tipbox}
\textbf{Method:} Scale each person's fraction up to a full job first --- A: $12/(3/5)=20$ days alone;
B: $6/(2/5)=15$ days alone. Then combine as usual: together $=(20\times15)/(20+15)=60/7$ days.

\subsection*{K4. Reduced Efficiency Due to a Mid-Job Slowdown}
\begin{tipbox}[Phrasing]
``A can finish a job in 15 days working normally. After 5 days, A starts working at half the usual
efficiency for the rest of the job. How many total days does A take?''
\end{tipbox}
\textbf{Method:} Split into two phases with two different rates --- normal rate first, halved rate
second; don't try to average the rates. First 5 days at $1/15$/day complete $1/3$; remaining $2/3$ at
half rate ($1/30$/day) takes $20$ more days $\Rightarrow$ \textbf{Total = 25 days}.

\subsection*{K5. Repeating Break / Rest Cycle}
\begin{tipbox}[Phrasing]
``A works for 5 days and rests on the 6th (repeating). A alone can finish a job in 24 days of actual
work. How many total days (including rest days) will the job take?''
\end{tipbox}
\textbf{Method:} Treat the rest day as adding to the calendar, not to the work. 4 full 6-day cycles
cover $4\times5=20$ working days in $4\times6=24$ calendar days; 4 working days remain, needing 4 more
calendar days $\Rightarrow$ \textbf{Total = 28 calendar days}. Always separate ``days worked'' from
``days elapsed'' for this variant.

\subsection*{K6. Machine Breaks Down and Is Repaired}
\begin{tipbox}[Phrasing]
``A machine can complete a job in 12 hours. It runs for 3 hours, then breaks down and is under repair
for 2 hours (doing no work), then resumes until the job is finished. How long does the whole job
take, start to finish?''
\end{tipbox}
\textbf{Method:} Split into phases exactly like a pipe/machine problem, but insert a zero-work phase
for the downtime. Work done in first 3 hours $=3/12=1/4$; remaining $3/4$ at the same rate needs $9$
more running hours $\Rightarrow$ \textbf{Total elapsed = 3 + 2 (repair) + 9 = 14 hours}.

\subsection*{K7. Wrong Time Assumption / Efficiency-Loss Question}
\begin{tipbox}[Phrasing]
``A estimated finishing a job in 10 days at a certain daily rate, but actually took 12 days at a
lower daily rate. By what percentage did A's actual daily efficiency fall short of the estimate?''
\end{tipbox}
\textbf{Method:} Efficiency is inversely proportional to time (\texttt{01-fundamentals.tex} \S6).
Estimated rate $\propto 1/10$, actual rate $\propto 1/12$. Percentage shortfall
$=(1/10-1/12)/(1/10)\times100=16.67\%$.

\subsection*{K8. Workforce Increases by a Percentage (Not Efficiency)}
\begin{tipbox}[Phrasing]
``20 workers can complete a job in 18 days. If the number of workers is increased by 20\%, in how
many days will the job be completed?''
\end{tipbox}
\textbf{Method:} $\text{Workers}\times\text{Days}=\text{Constant}$ (same worker-day logic as
\texttt{02-fast-tricks-and-shortcuts.tex} \S3, applied to headcount instead of a mid-job change). New
workers $=20\times1.2=24$; new days $=(20\times18)/24=15$ \textbf{days}. Keep this distinct from an
\emph{efficiency} increase --- here the per-person rate is unchanged, only headcount changes.

\subsection*{K9. Unit-Production Phrasing (Bottles, Units, Items per Hour)}
\begin{tipbox}[Phrasing]
``Machine A produces 200 bottles per hour, Machine B produces 150 bottles per hour. Working together,
how long will they take to produce 7000 bottles?''
\end{tipbox}
\textbf{Method:} Identical to a combined-work question --- treat ``bottles'' as the job. Combined
rate $=200+150=350$ bottles/hour; time $=7000/350=20$ \textbf{hours}. Don't be thrown by the absence
of the word ``work'' --- any per-unit-time output phrasing (bottles, pages, units, bricks) uses the
same rate-addition method.

\subsection*{K10. Mixed Skilled / Unskilled Workers}
\begin{tipbox}[Phrasing]
``3 skilled workers and 5 unskilled workers can complete a job in 6 days. A skilled worker is twice
as efficient as an unskilled worker. If 4 skilled workers alone were to do the job, how many days
would they take?''
\end{tipbox}
\textbf{Method:} Convert every worker to one common unit first. Let unskilled rate $=1$ unit/day
$\Rightarrow$ skilled rate $=2$ units/day. Combined rate $=3(2)+5(1)=11$ units/day; total work
$=11\times6=66$ units. 4 skilled workers' rate $=4\times2=8$ units/day $\Rightarrow$ time
$=66/8=8.25$ \textbf{days}.

\section*{Extra Trigger-Word Tricks}

These are quick mental reflexes to pair with the ten variants above --- the same ``see the word, do
the step'' style as \texttt{02-fast-tricks-and-shortcuts.tex}.

\begin{tipbox}[``...then leaves / then joins'']
Completed work $\to$ Remaining work. Never solve the whole job in one pass; split into phases at the
moment the workforce changes.
\end{tipbox}

\begin{tipbox}[``...twice / thrice / half as efficient'']
Convert straight to a rate ratio ($2{:}1$, $3{:}1$, $1{:}2$). Never convert to a day ratio first ---
that inversion step is where most errors happen.
\end{tipbox}

\begin{tipbox}[``...every Nth day / alternate days / repeating cycle'']
Think in whole cycles, not day-by-day. Compute work done per cycle, find how many full cycles fit,
then solve only the small leftover.
\end{tipbox}

\begin{tipbox}[``...ratio of times / ratio of days is given'']
Immediately flip it: time ratio and efficiency ratio are always inverses of each other.
\end{tipbox}

\begin{tipbox}[25\% / 40\% / 75\% (any clean percentage)]
Convert to a fraction before touching the arithmetic ($25\%\to1/4$, $40\%\to2/5$, $75\%\to3/4$).
Fractions combine faster than decimals in every one of these variants.
\end{tipbox}

\begin{tipbox}[All given times are whole numbers]
Default to the LCM Method from \texttt{02-fast-tricks-and-shortcuts.tex} \S1 --- it applies just as
well to K3, K9, and K10 as it does to the core combined-work questions.
\end{tipbox}

\begin{tipbox}[Pipes machines or downtime in the same question]
Keep fill/production as $+$ and empty/downtime as a $0$-rate phase, and always track \emph{elapsed
time} separately from \emph{working time} when a repair or rest period is involved (K5, K6).
\end{tipbox}

\section*{What This File Deliberately Skips}

Per the original review, these remain \textbf{not worth prepping for TCS} and are intentionally
excluded here as well: multi-page algebraic worker puzzles, hour-by-hour efficiency changes,
irregular A/B/C/D rotation schedules, infinite pipe on/off sequences, nonlinear efficiency models,
and minimum-time optimization/scheduling problems. These belong to CAT/SSC CGL-level prep, not TCS
NQT.

\section*{Updated Priority Order}

With this file added, the full priority order across all four booklets becomes:

\begin{enumerate}[leftmargin=1.4em]
  \item \textbf{A1} --- Two people working together (base formula)
  \item \textbf{B1} --- Efficiency ratio (``X times as efficient'')
  \item \textbf{C1/C2} --- Someone joins or leaves midway
  \item \textbf{K1/K2} --- Working for x days then leaving / together-then-one-leaves (same family
  as C1/C2)
  \item \textbf{D1} --- Men/Women/Children equivalence
  \item \textbf{E1/E2} --- Pipes and cisterns, including a pipe closing partway
  \item \textbf{F1} --- Alternate-day working
  \item \textbf{K5} --- Repeating break/rest cycle (same family as F1)
  \item \textbf{G1} --- Wages split by work share
  \item \textbf{H1} --- Machine/Printer/Robot work
  \item \textbf{K6} --- Machine breaks down and is repaired (same family as H1)
  \item \textbf{I1} --- Percentage of work completed
  \item \textbf{K3} --- Fraction of work completed by each person (same family as I1)
  \item \textbf{K9} --- Unit-production phrasing (bottles/pages/units)
  \item \textbf{J1/J2} --- Variable efficiency
  \item \textbf{K4} --- Reduced efficiency due to mid-job slowdown (same family as J1/J2)
  \item \textbf{K8} --- Workforce increases by a percentage
  \item \textbf{K10} --- Mixed skilled/unskilled workers
  \item \textbf{K7} --- Wrong time assumption / efficiency-loss question (lowest frequency, Digital
  only)
\end{enumerate}

\noindent At this point the coverage is effectively complete for TCS NQT. The highest-value next step
is pure repetition: solve 40--60 mixed questions across all four booklets until each pattern is
recognized on sight, since TCS rewards speed of recognition over depth of math.

\end{document}
